\documentclass[11pt,a4paper]{article}

\usepackage[utf8]{inputenc}
\usepackage[T1]{fontenc}
\usepackage[german]{babel}
\usepackage{amsmath,amssymb,amsthm}
\usepackage{graphicx}
\usepackage{hyperref}
\usepackage{geometry}
\geometry{margin=2.5cm}

\title{Eine zweiskalige Struktur der Faktorisierungskomplexität\
	\large Empirische Evidenz für die Rolle arithmetischer Glattheit}
\author{Thomas Hoffbauer}
\date{\today}

\begin{document}
	
	\maketitle
	
	\begin{abstract}
		Wir untersuchen empirisch die Laufzeit von Faktorisierungsalgorithmen auf zufälligen ganzen Zahlen.
		Die Daten zeigen deutlich, dass die Laufzeit nicht allein durch die Größe der Zahl bestimmt wird,
		sondern stark von der arithmetischen Struktur abhängt, insbesondere vom größten Primfaktor.
		Wir identifizieren eine zweiskalige Struktur: eine globale Skalierung durch $\log n$ und eine lokale Korrektur durch die Glattheit der Zahl.
		Dies führt zu einer natürlichen Klassifikation von Zahlen in unterschiedliche ``Rechenphasen''.
	\end{abstract}
	
	\section{Einleitung}
	
	Die Komplexität der Faktorisierung ganzer Zahlen ist ein zentrales Problem der algorithmischen Zahlentheorie.
	Klassische Resultate zeigen, dass moderne Verfahren eine asymptotische Laufzeit der Form
	
	\begin{equation}
		T(n) \sim \exp\left( \sqrt{\log n \cdot \log \log n} \right)
	\end{equation}
	
	aufweisen.
	
	Diese Beschreibung berücksichtigt jedoch nur die Größe der Zahl und ignoriert deren interne Struktur.
	In dieser Arbeit zeigen wir, dass diese Struktur entscheidend ist.
	
	\section{Empirische Beobachtungen}
	
	Wir analysieren Laufzeiten verschiedener Faktorisierungsverfahren auf zufällig generierten Zahlen unterschiedlicher Größe.
	
	Die wichtigsten Beobachtungen sind:
	
	\begin{itemize}
		\item Die Laufzeit zeigt deutliche Cluster.
		\item Primzahlen und semiprime Zahlen sind signifikant schwerer zu faktorisieren.
		\item Zahlen mit vielen kleinen Primfaktoren sind besonders leicht.
	\end{itemize}
	
	Diese Struktur bleibt auch bei fixer Stellenzahl erhalten.
	
	\section{Glattheit als strukturgebender Parameter}
	
	Wir definieren die Glattheit einer Zahl über ihren größten Primfaktor:
	
	\begin{equation}
		P_{\max}(n) = \max { p \mid p \text{ Primteiler von } n }
	\end{equation}
	
	Empirisch zeigt sich eine starke Korrelation:
	
	\begin{equation}
		\text{Runtime}(n) \sim \log(P_{\max}(n))
	\end{equation}
	
	Insbesondere gilt:
	
	\begin{itemize}
		\item Kleine $P_{\max}$ $\Rightarrow$ schnelle Faktorisierung
		\item Große $P_{\max}$ $\Rightarrow$ langsame Faktorisierung
	\end{itemize}
	
	\section{Zweiskaliges Modell}
	
	Wir schlagen folgendes Modell vor:
	
	\begin{equation}
		T(n) \sim \exp\left(
		\sqrt{\log n \cdot \log\log n}
		\cdot g(P_{\max}(n))
		\right)
	\end{equation}
	
	wobei $g$ eine monotone Funktion ist.
	
	Eine einfache Näherung ist:
	
	\begin{equation}
		T(n) \approx A \cdot \exp(B \cdot \sqrt{\log n \log\log n}) \cdot (P_{\max}(n))^\alpha
	\end{equation}
	
	\section{Phasenstruktur}
	
	Die Daten legen nahe, dass Zahlen in drei Klassen zerfallen:
	
	\begin{enumerate}
		\item Glatte Zahlen (kleine Primfaktoren)
		\item Gemischte Zahlen
		\item Primzahlartige Zahlen (große Primfaktoren)
	\end{enumerate}
	
	Diese Klassen bilden klar getrennte Cluster im Merkmalsraum.
	
	\section{Maschinelles Lernen}
	
	Ein einfacher Random-Forest-Klassifikator kann die Schwierigkeit einer Faktorisierung mit hoher Genauigkeit vorhersagen, basierend auf:
	
	\begin{itemize}
		\item modularen Restklassen
		\item Ziffernentropie
		\item kleinen Teilbarkeitstests
	\end{itemize}
	
	Dies zeigt, dass die Struktur der Zahl bereits vor der Faktorisierung zugänglich ist.
	
	\section{Interpretation}
	
	Die Ergebnisse legen nahe:
	
	\begin{itemize}
		\item Faktorisierung ist kein homogenes Problem
		\item sondern ein strukturabhängiger Prozess
		\item mit klaren ``Phasenübergängen''
	\end{itemize}
	
	Der größte Primfaktor wirkt als eine Art ``Barriere'' im Rechenprozess.
	
	\section{Schlussfolgerung}
	
	Wir zeigen, dass die klassische Komplexitätstheorie durch einen strukturellen Term ergänzt werden muss.
	Die Glattheit einer Zahl ist ein zentraler Parameter für die tatsächliche Rechenzeit.
	
	Dies eröffnet neue Möglichkeiten für:
	
	\begin{itemize}
		\item adaptive Faktorisierungsalgorithmen
		\item Vorhersagemodelle
		\item theoretische Verallgemeinerungen
	\end{itemize}
	
	\section*{Ausblick}
	
	Zukünftige Arbeiten sollten:
	
	\begin{itemize}
		\item die Funktion $g$ explizit bestimmen
		\item den Zusammenhang mit bekannten Verteilungen glatter Zahlen untersuchen
		\item die Phasenstruktur theoretisch begründen
	\end{itemize}
	
\end{document}
